% Осень26!A16; нумерация внутри столбца без пустых строк.
% Источник: лекция 6, §5 (строка 295)
\ticket{14}{Теорема о сохранении области}
\label{ticket:14}

\subsection*{\texorpdfstring{Возврат в \(\mathbb R^m\). Теорема о сохранении области}{Возврат в \textbackslash mathbb R\^{}m. Теорема о сохранении области}}

«Вернёмся в наше старое доброе \(\mathbb R^m\)» (лектор). Замечание о названии: «область» обычно означает открытое \textbf{связное} множество; в теореме связность не участвует --- только открытость.

\begin{lemma}

\emph{о почти локальной инъективности --- из предыдущей лекции, цитируется.} Пусть \(O\subset\mathbb R^m\) открыто, \(f\colon O\to\mathbb R^m\) дифференцируемо в точке \(x_0\in O\) и производный оператор \(f'(x_0)\) невырожден. Тогда существуют \(C>0\) и \(\delta>0\) такие, что \(\bar B(x_0,\delta)\subset O\) и \[\forall h,\ \|h\|\le\delta:\qquad \bigl\|f(x_0+h)-f(x_0)\bigr\|\ \ge\ C\,\|h\|.\] \textbf{Требует сверки.} Формулировка восстановлена по цитированию в лекции : лектор оговаривал, что берёт \textbf{замкнутый} шар (в лемме был открытый --- «такое цитирование ничего не искажает»), и уточнял на доске знаки/неравенства по ходу.

\end{lemma}

\begin{theorem}

\emph{о сохранении области.} Пусть \(O\subset\mathbb R^m\) открыто, \(f\colon O\to\mathbb R^m\), и в \textbf{каждой} точке \(x\in O\) существует производный оператор \(f'(x)\), задаваемый квадратной матрицей, и он \textbf{невырожден}: \[\det f'(x)\ne 0\quad(\Leftrightarrow\ f'(x)\ \text{обратим})\qquad\forall x\in O.\] Тогда \(f(O)\) открыто в \(\mathbb R^m\).

\end{theorem}

\begin{proof}

Пусть \(y_0\in f(O)\), т.е. \(y_0=f(x_0)\), \(x_0\in O\). Докажем, что \(y_0\) --- внутренняя точка \(f(O)\).

\textbf{(1) Сфера не схлопывается в точку.} По лемме о почти локальной инъективности найдутся \(C>0\), \(\delta>0\): \(\bar B(x_0,\delta)\subset O\) и \(\|f(x)-f(x_0)\|\ge C\|x-x_0\|\) на \(\bar B(x_0,\delta)\). В частности, на сфере \(S(x_0,\delta):=\{x:\|x-x_0\|=\delta\}\) выполнено \(f(x)\ne f(x_0)=y_0\) (и даже \(\|f(x)-y_0\|\ge C\delta\)).

\textbf{(2) Расстояние до образа сферы положительно.} Положим \[R:=\tfrac12\,\mathrm{dist}\bigl(y_0,\,f(S(x_0,\delta))\bigr):=\tfrac12\inf\{\|y_0-u\|:\ u\in f(S(x_0,\delta))\}.\] Утверждение: \(R>0\). Обоснование (существенно использует компактность!):

\begin{itemize}
\tightlist
\item
  сфера \(S(x_0,\delta)\) замкнута (уравнение сферы сохраняется предельным переходом) и ограничена, т.е. \textbf{компактна} в \(\mathbb R^m\);
\item
  \(f\) непрерывна (дифференцируемость влечёт непрерывность), поэтому образ \(f(S(x_0,\delta))\) --- компакт;
\item
  функция \(u\mapsto\|y_0-u\|\) непрерывна на компакте и всюду \(>0\) (точки образа не совпадают с \(y_0\) по п. (1)), поэтому по теореме Вейерштрасса она \textbf{достигает} минимума, и минимум положителен.
\end{itemize}

Контраст-иллюстрация лектора: \(\inf\{x:\ x\in(0,1)\}=0\) --- для положительной функции на \textbf{некомпактном} множестве инфимум может быть нулём (сколь угодно близкие приближения). Здесь спасает компактность: инфимум есть минимум.

\textbf{(3) Цель: показать \(B(y_0,R)\subset f(O)\).} Зафиксируем произвольное \(y\) с \(\|y-y_0\|<R\) и будем искать \(x\in B(x_0,\delta)\) с \(f(x)=y\). Рассмотрим вспомогательную (уже числовую!) функцию \[g(x):=\|f(x)-y\|^2=\sum_{j=1}^{m}\bigl(f_j(x)-y_j\bigr)^2,\qquad x\in\bar B(x_0,\delta)\] (квадрат нормы вместо нормы --- «технические причины»: исчезают квадратные корни). Мечта: \(g\) где-нибудь обнуляется --- это и есть \(f(x)=y\).

\textbf{(4) Два наблюдения о значениях \(g\).}

\begin{itemize}
\item
  \begin{enumerate}
  \def\labelenumi{(\roman{enumi})}
  \tightlist
  \item
    В центре: \(g(x_0)=\|f(x_0)-y\|^2=\|y_0-y\|^2<R^2\).
  \end{enumerate}
\item
  \begin{enumerate}
  \def\labelenumi{(\roman{enumi})}
  \setcounter{enumi}{1}
  \tightlist
  \item
    На сфере \(S(x_0,\delta)\): по неравенству треугольника и определению \(R\) (для \(x\in S\): \(\|f(x)-y_0\|\ge\mathrm{dist}(y_0,f(S))=2R\)): \[\|f(x)-y\|\ \ge\ \|f(x)-y_0\|-\|y_0-y\|\ >\ 2R-R=R,\] т.е. \(g(x)\ge R^2\) на сфере. \emph{(Геометрическое объяснение лектора: всякий отрезок, идущий из точки образа сферы --- отстоящей от \(y_0\) не менее чем на \(2R\) --- внутрь шара \(B(y_0,R)\), имеет длину больше \(R\): «строго по радиусу --- ровно \(R\), а с отклонением --- по неравенству треугольника больше».)}
  \end{enumerate}
\end{itemize}

\textbf{(5) Минимум внутри.} \(g\) непрерывна на компакте \(\bar B(x_0,\delta)\) \ensuremath{\Longrightarrow} по теореме Вейерштрасса минимум достигается в некоторой точке \(x^*\). Из (i)--(ii): на границе (сфере) все значения \(g\ge R^2>g(x_0)\), значит минимум достигается \textbf{строго внутри}: \(x^*\in B(x_0,\delta)\subset O\).

\textbf{(6) Теорема Ферма + невырожденность.} Внутренняя точка минимума --- критическая: по теореме Ферма все частные производные \(g\) в \(x^*\) нулевые. Дифференцируя (сократив двойку): \[0=\frac{\partial g}{\partial x_k}(x^*)=\sum_{j=1}^{m}\bigl(f_j(x^*)-y_j\bigr)\,\frac{\partial f_j}{\partial x_k}(x^*)\qquad k=1,\dots,m.\] Это --- \textbf{однородная линейная система}: в «роли неизвестных» --- скобки \((f_j(x^*)-y_j)\), в «роли коэффициентов» --- частные производные (вычисленные в искомой точке \(x^*\)); матрица системы --- транспонированная матрица Якоби \(f'(x^*)\). Операторная запись: \[\bigl(f'(x^*)\bigr)^{\mathsf T}\,\bigl(f(x^*)-y\bigr)=0.\] Невырожденный оператор подействовал на вектор и дал нуль \ensuremath{\Longrightarrow} вектор нулевой: \[f(x^*)-y=0\quad\Longrightarrow\quad y=f(x^*)\in f(O).\] (Отсюда, кстати, минимум \(g\) равен нулю --- но важна не величина, а сам факт разрешимости.)

\textbf{(7) Итак,} для произвольного \(y\in B(y_0,R)\) нашли прообраз; значит \(B(y_0,R)\subset f(O)\), и \(y_0\in f(O)\) --- внутренняя точка. Так как \(y_0\) произвольна, \(f(O)\) открыто.

\end{proof}

\textbf{Наглядная картинка лектора} («резиночка»): непрерывное отображение деформирует лист-резинку; вопрос --- может ли край (образ сферы) «улететь» вбок, оставив рядом с \(y_0\) дыру? Доказательство показывает, что нет: шар \(B(y_0,R)\) целиком заполнен образами.
